\section{The evolution operator and its coefficients}\label{sec:omega}
We now introduce the evolution operator from a field-theoretic perspective. The construction follows
the operator-based approach of Ref.~\cite{Kovner:2018rbf}, in which the evolution operator is a
unitary transformation that both generates the soft gluon wave function and diagonalises the
interacting Yang--Mills Hamiltonian order by order in perturbation theory.

The eigenstates of the interacting Hamiltonian are not the free Fock states. The interaction
Hamiltonian mixes states with different numbers of gluons, so the physical states are dressed states,
i.e.\ superpositions of multi-gluon Fock components. This dressing is implemented systematically by
a unitary operator. With $H=H_{0}+H_{\rm int}$ the light-cone Hamiltonian of Sec.~\ref{subsec:LCQCD},
the interacting eigenstates are related to the free Fock states by
\begin{align}\label{eq:dressing}
\ket{\Psi_{n}}=\Omega\,\ket{n}\,,\qquad \Omega^{\dagger}\Omega=\Omega\,\Omega^{\dagger}=1\,,
\end{align}
where $\ket{n}$ is a free Fock state of the soft sector, $\ket{\Psi_{n}}$ the corresponding dressed
eigenstate of $H$, and $\Omega$ the evolution operator. In particular, the dressed soft vacuum of
Eq.~\eqref{eq:Wfunc} is $\ket{\Psi}=\Omega\ket{0}$. Like any operator on the soft Fock space, $\Omega$ can
be written as a normal-ordered series in the soft gluon creation and annihilation operators. In the
eikonal approximation the valence sector enters the Hamiltonian only through $\rho$, so the
coefficients of this series are functionals of $\rho$; they are operators on the valence space, and
their ordering matters. Following Ref.~\cite{Radhakrishnan:2026bxy},
\begin{align}\label{eq:OmegaExp}
\Omega=\ &\mf{N}
+\int\frac{d^3p}{(2\pi)^3}\Big[\mf{A}^b_{1j}(p)\,a^b_j(p)+\mf{A}^b_{2j}(p)\,a^{\dagger b}_j(p)\Big]
\notag\\
&+\int\frac{d^3p}{(2\pi)^3}\frac{d^3q}{(2\pi)^3}\Big[\mf{B}^{cb}_{1kj}(p,q)\,a^b_j(p)\,a^c_k(q)
+\mf{B}^{cb}_{2kj}(p,q)\,a^{\dagger b}_j(p)\,a^{\dagger c}_k(q)
\notag\\
&\hspace{3.4cm}+\mf{B}^{cb}_{3kj}(p,q)\,a^{\dagger c}_k(q)\,a^b_j(p)\Big]
\notag\\
&+\int\frac{d^3p}{(2\pi)^3}\frac{d^3q}{(2\pi)^3}\Big[\mf{C}^{dcb}_{1lkj}(p,q)\,a^{\dagger d}_l(q)\,a^{\dagger c}_k(p-q)\,a^b_j(p)
\notag\\
&\hspace{3.4cm}+\mf{C}^{dcb}_{2lkj}(p,q)\,a^{\dagger b}_j(p)\,a^c_k(p-q)\,a^d_l(q)\Big]+\dots,
\end{align}
where the dots denote terms with more soft gluon operators, which do not contribute to the single
inclusive cross section at order $g^{4}$. The longitudinal integrals run over the soft window,
$\Lambda<p^{+},q^{+}<\vee$. The colour and polarisation indices are tied to the momenta as follows: in
$\mf{B}^{cb}_{kj}(p,q)$ the pair $(b,j)$ belongs to $p$ and $(c,k)$ to $q$; in
$\mf{C}^{dcb}_{lkj}(p,q)$, $(b,j)$ belongs to $p$, $(d,l)$ to $q$ and $(c,k)$ to $p-q$. In the
three-operator terms we have used momentum conservation in the soft sector, which holds for the vertices
that do not involve $\rho$. This is sufficient here because $\mf{C}_{1,2}$ are needed only at order
$g$, where they are independent of $\rho$. Since two creation (or two annihilation) operators commute,
only the parts of the coefficients symmetric under their exchange are defined, and we impose
\begin{align}\label{eq:symm}
\mf B^{cb}_{1kj}(p,q)=\mf B^{bc}_{1jk}(q,p)\,,\qquad \mf B^{cb}_{2kj}(p,q)=\mf B^{bc}_{2jk}(q,p)\,,\qquad
\mf C^{dcb}_{1lkj}(p,q)=\mf C^{cdb}_{1klj}(p,p-q)\,,
\end{align}
and the same relation for $\mf C_{2}$.

The coefficients have definite orders in $g$, and some of them contain two orders. To the accuracy
needed here,
\begin{align}\label{eq:orders}
\mf{N}=1+\mf{N}^{(2)}\,,\qquad
\mf{A}^{b}_{1j}=\mf{A}^{(1)b}_{1j}+\mf{A}^{(3)b}_{1j}\,,\qquad
\mf{A}^{b}_{2j}=\mf{A}^{(1)b}_{2j}+\mf{A}^{(3)b}_{2j}\,,\qquad
\mf{B}=\mathcal{O}(g^{2})\,,\qquad \mf{C}=\mathcal{O}(g)\,,
\end{align}
where the superscript gives the order in $g$, and $\mf{N}$ is Hermitian. We need $\Omega$ to order
$g^{3}$, but not all of its $g^{3}$ terms: we keep only those terms that produce at most two soft
gluons when $\Omega$ acts on the vacuum or on a one-gluon state, since states with more gluons do not
contribute to single inclusive production. The coefficients are fixed by two complementary sets of
conditions:
\begin{itemize}
\item unitarity, $\Omega^{\dagger}\Omega=1$, imposed order by order in $g$;
\item the requirement that $\Omega$ acting on the free Fock states reproduces the dressed states
computed directly in LCPT, $\Omega\ket{n}=\ket{\Psi_{n}}$, for the vacuum and the one-gluon states.
\end{itemize}

\subsection{Unitarity constraints}\label{subsec:unitarity}
Normal ordering $\Omega^{\dagger}\Omega$ with Eq.~\eqref{eq:aadagger} and keeping the terms required
for the coefficients at the orders of Eq.~\eqref{eq:orders}, we find
\begin{align}\label{eq:unitarity}
1=\Omega^{\dagger}\Omega=\ &\mf{N}^{2}+\int\frac{d^{3}p}{(2\pi)^{3}}\,\mf{A}^{(1)\dagger b}_{2j}(p)\,\mf{A}^{(1)b}_{2j}(p)\notag\\
&+\int\frac{d^{3}p}{(2\pi)^{3}}\,a^{\dagger b}_{j}(p)\bigg[\mf{A}^{(1)b}_{2j}(p)+\mf{A}^{(1)\dagger b}_{1j}(p)
+\mf{A}^{(3)b}_{2j}(p)+\mf{A}^{(3)\dagger b}_{1j}(p)
+\mf{N}^{(2)}\mf{A}^{(1)b}_{2j}(p)+\mf{A}^{(1)\dagger b}_{1j}(p)\,\mf{N}^{(2)}\notag\\
&\hspace{3.2cm}+\int\frac{d^{3}q}{(2\pi)^{3}}\Big(2\,\mf{A}^{(1)\dagger c}_{2k}(q)\,\mf{B}^{cb}_{2kj}(p,q)
+\mf{B}^{\dagger cb}_{3kj}(p,q)\,\mf{A}^{(1)c}_{2k}(q)\notag\\
&\hspace{3.2cm}+2\,\mf{C}^{\dagger dcb}_{1lkj}(p,q)\,\mf{B}^{cd}_{2kl}(q,p-q)\Big)\bigg]\notag\\
&+\int\frac{d^{3}p}{(2\pi)^{3}}\frac{d^{3}q}{(2\pi)^{3}}\,a^{\dagger b}_{j}(p)\,a^{\dagger c}_{k}(q)\bigg[\mf{B}^{cb}_{2kj}(p,q)+\mf{B}^{\dagger cb}_{1kj}(p,q)
+\frac{1}{2}\,\mf{A}^{(1)\dagger b}_{1j}(p)\,\mf{A}^{(1)c}_{2k}(q)\notag\\
&\hspace{3.2cm}+\frac{1}{2}\,\mf{A}^{(1)\dagger c}_{1k}(q)\,\mf{A}^{(1)b}_{2j}(p)
+\mf{C}^{\dagger cbd}_{2kjl}(p+q,q)\,\mf{A}^{(1)d}_{2l}(p+q)\bigg]\notag\\
&+\int\frac{d^{3}p}{(2\pi)^{3}}\frac{d^{3}q}{(2\pi)^{3}}\,a^{\dagger b}_{j}(p)\,a^{c}_{k}(q)\bigg[\mf{B}^{\dagger cb}_{3kj}(p,q)+\mf{B}^{bc}_{3jk}(q,p)
+\mf{A}^{(1)\dagger b}_{1j}(p)\,\mf{A}^{(1)c}_{1k}(q)+\mf{A}^{(1)\dagger c}_{2k}(q)\,\mf{A}^{(1)b}_{2j}(p)\notag\\
&\hspace{3.2cm}+2\,\mf{A}^{(1)\dagger d}_{2l}(q-p)\,\mf{C}^{bdc}_{1jlk}(q,p)
+2\,\mf{C}^{\dagger cdb}_{1klj}(p,q)\,\mf{A}^{(1)d}_{2l}(p-q)\bigg]\notag\\
&+\int\frac{d^{3}p}{(2\pi)^{3}}\frac{d^{3}q}{(2\pi)^{3}}\,a^{\dagger d}_{l}(q)\,a^{\dagger c}_{k}(p-q)\,a^{b}_{j}(p)\Big[\mf{C}^{\dagger dcb}_{2lkj}(p,q)+\mf{C}^{dcb}_{1lkj}(p,q)\Big]
+{\rm h.c.}+\dots,
\end{align}
where ``h.c.'' denotes the Hermitian conjugates of the terms with only creation operators. Each
bracket must vanish order by order in $g$. This gives five relations:
\begin{align}
\mf{N}&=1-\frac{1}{2}\int\frac{d^{3}p}{(2\pi)^{3}}\,\mf{A}^{(1)\dagger b}_{2j}(p)\,\mf{A}^{(1)b}_{2j}(p)
=1-\frac{1}{2}\int\frac{d^{3}p}{(2\pi)^{3}}\,\mf{A}^{(1)b}_{1j}(p)\,\mf{A}^{(1)\dagger b}_{1j}(p)\,,\label{eq:Nrel}\\
\mf{A}^{(1)\dagger b}_{1j}(p)&=-\mf{A}^{(1)b}_{2j}(p)\,,\label{eq:A1rel}\\
\mf{A}^{(3)\dagger b}_{1j}(p)&=-\mf{A}^{(3)b}_{2j}(p)+\big[\mf{A}^{(1)b}_{2j}(p),\mf{N}^{(2)}\big]
+\int\frac{d^{3}q}{(2\pi)^{3}}\Big(2\,\mf{A}^{(1)c}_{1k}(q)\,\mf{B}^{cb}_{2kj}(p,q)
-\mf{B}^{\dagger cb}_{3kj}(p,q)\,\mf{A}^{(1)c}_{2k}(q)\notag\\
&\hspace{6.1cm}-2\,\mf{C}^{\dagger dcb}_{1lkj}(p,q)\,\mf{B}^{cd}_{2kl}(q,p-q)\Big)\,,\label{eq:A3rel}\\
\mf{B}^{\dagger cb}_{1kj}(p,q)&=\frac{1}{2}\Big[-2\,\mf{B}^{cb}_{2kj}(p,q)+\mf{A}^{(1)b}_{2j}(p)\,\mf{A}^{(1)c}_{2k}(q)+\mf{A}^{(1)c}_{2k}(q)\,\mf{A}^{(1)b}_{2j}(p)
+2\,\mf{C}^{cbd}_{1kjl}(p+q,q)\,\mf{A}^{(1)d}_{2l}(p+q)\Big]\,,\label{eq:B1rel}\\
\mf{C}^{\dagger dcb}_{2lkj}(p,q)&=-\mf{C}^{dcb}_{1lkj}(p,q)\,.\label{eq:C2rel}
\end{align}
In Eq.~\eqref{eq:A3rel} we have used Eq.~\eqref{eq:A1rel} and the hermiticity of $\mf{N}$, so that
$-\mf{N}^{(2)}\mf{A}^{(1)}_{2}-\mf{A}^{(1)\dagger}_{1}\mf{N}^{(2)}=[\mf{A}^{(1)}_{2},\mf{N}^{(2)}]$.
This commutator does not vanish. Both $\mf{N}^{(2)}$ and $\mf{A}^{(1)}_{2}$ are operators on the valence
space, $\mf{N}^{(2)}$ quadratic and $\mf{A}^{(1)}_{2}$ linear in $\rho$. With the explicit coefficients
obtained in Sec.~\ref{sec:vacin},
\begin{align}\label{eq:A1N2explicit}
\mf{A}^{(1)b}_{2j}(p)=-\frac{\sqrt{2}\,g\,\v p^{j}}{\sqrt{p^{+}}\,\v p^{2}}\,\rho^{b}(-\v p)\,,\qquad
\mf{N}^{(2)}=-\frac{g^{2}}{(2\pi)^{3}}\,\ln\Big(\frac{\vee}{\Lambda}\Big)\int d^{2}\v q\,
\frac{\rho^{a}(-\v q)\,\rho^{a}(\v q)}{\v q^{2}}\,,
\end{align}
and the algebra~\eqref{eq:rhoalgebra}, we find
\begin{align}\label{eq:NAcomm}
\big[\mf{A}^{(1)b}_{2j}(p),\mf{N}^{(2)}\big]=-\frac{i\,g^{3}}{8\sqrt{2}\,\pi^{3}}\,\ln\Big(\frac{\vee}{\Lambda}\Big)\,
\frac{\v p^{j}}{\sqrt{p^{+}}\,\v p^{2}}\,f^{abc}\int d^{2}\v q\,
\bigg[\frac{1}{\v q^{2}}-\frac{1}{(\v q-\v p)^{2}}\bigg]\big\{\rho^{a}(-\v q),\rho^{c}(\v q-\v p)\big\}\,.
\end{align}
To write it in this symmetric form we have used that $f^{abc}\{\rho^{a}(-\v q),\rho^{c}(\v q-\v p)\}$ is odd
under $\v q\to\v p-\v q$. The weight in square brackets is odd as well, so the integrand is even and the
commutator is non-zero for $\v p\neq0$. It would vanish only for $\v p=0$, where $\rho^{b}(-\v p)$
becomes the total colour charge. The commutator is of order $g^{3}\rho^{2}\ln(\vee/\Lambda)$, the same
order as the two-$\rho$ part of $\mf{A}^{(3)}_{1}$, and we keep it throughout.

With these relations, $\Omega$ becomes
\begin{align}\label{eq:OmegaConstrained}
\Omega=\ &\mf{N}
+\int\frac{d^3p}{(2\pi)^3}\Big[\mf{A}^{(1)b}_{1j}(p)\,a^{b}_{j}(p)-\mf{A}^{(1)\dagger b}_{1j}(p)\,a^{\dagger b}_{j}(p)
+\mf{A}^{(3)b}_{1j}(p)\,a^{b}_{j}(p)+\mf{A}^{(3)b}_{2j}(p)\,a^{\dagger b}_{j}(p)\Big]\notag\\
&+\int\frac{d^3p}{(2\pi)^3}\frac{d^3q}{(2\pi)^3}\Big[\mf{B}^{cb}_{1kj}(p,q)\,a^{b}_{j}(p)\,a^{c}_{k}(q)
+\mf{B}^{cb}_{2kj}(p,q)\,a^{\dagger b}_{j}(p)\,a^{\dagger c}_{k}(q)
+\mf{B}^{cb}_{3kj}(p,q)\,a^{\dagger c}_{k}(q)\,a^{b}_{j}(p)\Big]\notag\\
&+\int\frac{d^3p}{(2\pi)^3}\frac{d^3q}{(2\pi)^3}\Big[\mf{C}^{dcb}_{1lkj}(p,q)\,a^{\dagger d}_{l}(q)\,a^{\dagger c}_{k}(p-q)\,a^{b}_{j}(p)
-\mf{C}^{\dagger dcb}_{1lkj}(p,q)\,a^{\dagger b}_{j}(p)\,a^{c}_{k}(p-q)\,a^{d}_{l}(q)\Big]+\dots,
\end{align}
where $\mf{A}^{(1)}_{2}$ and $\mf{C}_{2}$ have been eliminated with Eqs.~\eqref{eq:A1rel} and~\eqref{eq:C2rel}.
The remaining coefficients are fixed as follows. The dressed vacuum (Sec.~\ref{sec:vacin}) gives $\mf{N}$,
$\mf{A}^{(1)}_{2}=-\mf{A}^{(1)\dagger}_{1}$, $\mf{A}^{(3)}_{2}$ and $\mf{B}_{2}$. The dressed one-gluon state
gives $\mf{B}_{3}$ and $\mf{C}_{1}$. $\mf{B}_{1}$ then follows from Eq.~\eqref{eq:B1rel} and $\mf{A}^{(3)}_{1}$
from Eq.~\eqref{eq:A3rel}, including the commutator~\eqref{eq:NAcomm}. The dressed one-gluon state also
contains $\mf{A}_{1}$ directly, so Eqs.~\eqref{eq:Nrel}, \eqref{eq:A1rel} and~\eqref{eq:A3rel} serve as
checks of the LCPT results.

\subsection{Transverse coordinate representation}\label{subsec:omegacoord}
With the Fourier conventions of Eq.~\eqref{eq:FTconventions}, applied to the coefficients as c-number
functions of the transverse momenta, $\Omega$ reads
\begin{align}\label{eq:OmegaCoord}
\Omega=\ &\mf{N}
+\int\frac{dp^+}{2\pi}\int d^2x\,\Big[\mf{A}^{(1)b}_{1j}(p^{+},\v x)\,a^{b}_{j}(p^{+},\v x)
-\mf{A}^{(1)\dagger b}_{1j}(p^{+},\v x)\,a^{\dagger b}_{j}(p^{+},\v x)\notag\\
&\hspace{3.2cm}+\mf{A}^{(3)b}_{1j}(p^{+},\v x)\,a^{b}_{j}(p^{+},\v x)
+\mf{A}^{(3)b}_{2j}(p^{+},-\v x)\,a^{\dagger b}_{j}(p^{+},\v x)\Big]\notag\\
&+\int\frac{dp^+}{2\pi}\frac{dq^+}{2\pi}\int d^2x\,d^2y\,\Big[\mf{B}^{cb}_{1kj}(p^{+},\v x;q^{+},\v y)\,a^{b}_{j}(p^{+},\v x)\,a^{c}_{k}(q^{+},\v y)\notag\\
&\hspace{3.2cm}+\mf{B}^{cb}_{2kj}(p^{+},-\v x;q^{+},-\v y)\,a^{\dagger b}_{j}(p^{+},\v x)\,a^{\dagger c}_{k}(q^{+},\v y)\notag\\
&\hspace{3.2cm}+\mf{B}^{cb}_{3kj}(p^{+},\v x;q^{+},-\v y)\,a^{\dagger c}_{k}(q^{+},\v y)\,a^{b}_{j}(p^{+},\v x)\Big]\notag\\
&+\int\frac{dp^+}{2\pi}\frac{dq^+}{2\pi}\int d^2x\,d^2y\,d^2z\,\Big[\mf{C}^{dcb}_{1lkj}(p^{+},\v z-\v y;q^{+},\v y-\v x)\,
a^{\dagger d}_{l}(q^{+},\v x)\,a^{\dagger c}_{k}(p^{+}-q^{+},\v y)\,a^{b}_{j}(p^{+},\v z)\notag\\
&\hspace{3.2cm}-\mf{C}^{\dagger dcb}_{1lkj}(p^{+},\v z-\v y;q^{+},\v y-\v x)\,
a^{\dagger b}_{j}(p^{+},\v z)\,a^{c}_{k}(p^{+}-q^{+},\v y)\,a^{d}_{l}(q^{+},\v x)\Big]+\dots
\end{align}
The unitarity relations \eqref{eq:A1rel}--\eqref{eq:C2rel} become
\begin{align}
\mf{A}^{(1)\dagger b}_{1j}(p^{+},\v z)=\ &-\mf{A}^{(1)b}_{2j}(p^{+},-\v z)\,,\label{eq:A1relx}\\
\mf{A}^{(3)\dagger b}_{1j}(p^{+},\v x)=\ &-\mf{A}^{(3)b}_{2j}(p^{+},-\v x)+\big[\mf{A}^{(1)b}_{2j}(p^{+},-\v x),\mf{N}^{(2)}\big]\notag\\
&+\int\frac{dq^{+}}{2\pi}\int d^{2}y\,\Big[2\,\mf{A}^{(1)c}_{1k}(q^{+},\v y)\,\mf{B}^{cb}_{2kj}(p^{+},-\v x;q^{+},-\v y)\notag\\
&\hspace{2.4cm}-\mf{B}^{\dagger cb}_{3kj}(p^{+},\v x;q^{+},-\v y)\,\mf{A}^{(1)c}_{2k}(q^{+},-\v y)\notag\\
&\hspace{2.4cm}-2\int d^{2}z\,\mf{C}^{\dagger dcb}_{1lkj}(p^{+},\v x-\v z;q^{+},\v z-\v y)\,\mf{B}^{dc}_{2lk}(p^{+}-q^{+},-\v z;q^{+},-\v y)\Big]\,,\label{eq:A3relx}\\
\mf{B}^{\dagger cb}_{1kj}(p^{+},\v x;q^{+},\v y)=\ &\frac{1}{2}\Big[-2\,\mf{B}^{cb}_{2kj}(p^{+},-\v x;q^{+},-\v y)\notag\\
&\hspace{0.6cm}+\mf{A}^{(1)b}_{2j}(p^{+},-\v x)\,\mf{A}^{(1)c}_{2k}(q^{+},-\v y)+\mf{A}^{(1)c}_{2k}(q^{+},-\v y)\,\mf{A}^{(1)b}_{2j}(p^{+},-\v x)\notag\\
&\hspace{0.6cm}+2\int d^{2}z\,\mf{C}^{cbd}_{1kjl}(p^{+}+q^{+},\v z-\v x;q^{+},\v x-\v y)\,\mf{A}^{(1)d}_{2l}(p^{+}+q^{+},-\v z)\Big]\,,\label{eq:B1relx}\\
\mf{C}^{\dagger dcb}_{2lkj}(p^{+},\v z;q^{+},\v x)=\ &-\mf{C}^{dcb}_{1lkj}(p^{+},-\v z;q^{+},-\v x)\,.\label{eq:C2relx}
\end{align}
The commutator in Eq.~\eqref{eq:A3relx} is the transform of Eq.~\eqref{eq:NAcomm}. It can also be obtained
directly from the coordinate forms of $\mf{A}^{(1)}_{2}$ and $\mf{N}^{(2)}$, Eqs.~\eqref{lowav1}
and~\eqref{eq:Npos}, with
$[\rho^{a}(\v x),\rho^{b}(\v y)]=if^{abc}\rho^{c}(\v x)\,\delta^{(2)}(\v x-\v y)$:
\begin{align}\label{eq:NAcommx}
\big[\mf{A}^{(1)b}_{2j}(p^{+},-\v x),\mf{N}^{(2)}\big]=\ &\frac{g^{3}}{16\sqrt{2}\,\pi^{4}\sqrt{p^{+}}}\,
\ln\Big(\frac{\vee}{\Lambda}\Big)\,f^{bcd}\int_{\v u,\v y,\v w}
\bigg[\frac{(y-x)^{j}}{(y-x)^{2}}-\frac{(u-x)^{j}}{(u-x)^{2}}\bigg]\notag\\
&\times\frac{(u-w)\cdot(y-w)}{(u-w)^{2}\,(y-w)^{2}}\,\big\{\rho^{c}(\v u),\rho^{d}(\v y)\big\}\,.
\end{align}
The square bracket and $f^{bcd}\{\rho^{c}(\v u),\rho^{d}(\v y)\}$ are both odd under
$(\v u,c)\leftrightarrow(\v y,d)$, so the integrand is even and does not vanish. The two charges sit at
different transverse points, linked by the non-local kernel of $\mf{N}^{(2)}$. For a local kernel,
$\propto\delta^{(2)}(\v u-\v y)$, the commutator would vanish.

\subsection{Action on the soft vacuum and on a one-gluon state}\label{subsec:omegastates}
In momentum space the soft gluon states are
\begin{align}\label{eq:statesmom}
\ket{g^{b}_{j}(p)}=\frac{a^{\dagger b}_{j}(p)}{(2\pi)^{3/2}}\ket{0}\,,\qquad
\ket{g^{b}_{j}(p)\,g^{c}_{k}(q)}=\frac{a^{\dagger b}_{j}(p)\,a^{\dagger c}_{k}(q)}{(2\pi)^{3}}\ket{0}\,,
\end{align}
as follows from Eqs.~\eqref{eq:FTconventions} and~\eqref{eq:states}. Acting with Eq.~\eqref{eq:OmegaConstrained} on the soft
vacuum, only the terms without annihilation operators survive ($\mf{N}$ acts on the valence space only):
\begin{align}\label{eq:Omegavac}
\Omega\ket{0} =\ &\mf{N}\ket{0}+(2\pi)^{3/2}\int\frac{d^3p}{(2\pi)^3}\,\mf{A}^{b}_{2j}(p)\ket{g^{b}_{j}(p)}\notag\\
&+(2\pi)^3\int\frac{d^3p}{(2\pi)^3}\frac{d^3q}{(2\pi)^3}\,\mf{B}^{cb}_{2kj}(p,q)\ket{g^{b}_{j}(p)\,g^{c}_{k}(q)}+\dots,
\end{align}
where $\mf{A}_{2}=\mf{A}^{(1)}_{2}+\mf{A}^{(3)}_{2}=\mf{A}^{(3)}_{2}-\mf{A}^{(1)\dagger}_{1}$. In transverse
coordinate space,
\begin{align}\label{Omegaonvacpos}
\Omega\ket{0} =\ &\mf{N}\ket{0}
+\frac{1}{(2\pi)^{1/2}}\int dp^{+}\int d^{2}x\,
\Big[\mf{A}^{(3)b}_{2j}(p^{+},-\v x)-\mf{A}^{(1)\dagger b}_{1j}(p^{+},\v x)\Big]\ket{g^{b}_{j}(p^{+},\v x)}\notag\\
&+\frac{1}{2\pi}\int dp^{+}\,dq^{+}\int d^{2}z\,d^{2}x\,\mf{B}^{cb}_{2kj}(p^{+},-\v z;q^{+},-\v x)\,
\ket{g^{b}_{j}(p^{+},\v z)\,g^{c}_{k}(q^{+},\v x)}+\dots
\end{align}
In the two-gluon term we write $p^{+}=\xi k^{+}$ and $q^{+}=(1-\xi)k^{+}$, so that
$dp^{+}\,dq^{+}=k^{+}\,dk^{+}\,d\xi$:
\begin{align}\label{Omegaonvacposzeta}
\Omega\ket{0} =\ &\mf{N}\ket{0}
+\frac{1}{(2\pi)^{1/2}}\int dp^{+}\int d^{2}x\,
\Big[\mf{A}^{(3)b}_{2j}(p^{+},-\v x)-\mf{A}^{(1)\dagger b}_{1j}(p^{+},\v x)\Big]\ket{g^{b}_{j}(p^{+},\v x)}\notag\\
&+\frac{1}{2\pi}\int dk^{+}\,k^{+}\int d\xi\int d^{2}z\,d^{2}x\,
\mf{B}^{cb}_{2kj}\big(\xi k^{+},-\v z;(1-\xi)k^{+},-\v x\big)\,
\ket{g^{b}_{j}(\xi k^{+},\v z)\,g^{c}_{k}((1-\xi)k^{+},\v x)}+\dots
\end{align}
Acting on a one-gluon state $\ket{g^{a}_{i}(k)}$, one obtains
\begin{align}\label{eq:Omegag}
\Omega\ket{g^{a}_{i}(k)} =\ &\frac{1}{(2\pi)^{3/2}}\,\mf{A}^{a}_{1i}(k)\ket{0}+\mf{N}\ket{g^{a}_{i}(k)}
+\int\frac{d^3p}{(2\pi)^3}\,\mf{B}^{ba}_{3ji}(k,p)\ket{g^{b}_{j}(p)}\notag\\
&+(2\pi)^{3/2}\int\frac{d^3p}{(2\pi)^3}\,\mf{A}^{b}_{2j}(p)\ket{g^{a}_{i}(k)\,g^{b}_{j}(p)}\notag\\
&+(2\pi)^{3/2}\int\frac{d^3p}{(2\pi)^3}\,\mf{C}^{bca}_{1jki}(k,p)\ket{g^{b}_{j}(p)\,g^{c}_{k}(k-p)}+\dots,
\end{align}
with $\mf{A}_{1}=\mf{A}^{(1)}_{1}+\mf{A}^{(3)}_{1}$. In transverse coordinate space,
\begin{align}\label{Omegaongpos}
\Omega\ket{g^{a}_{i}(k^{+},\v z)} =\ &\frac{1}{(2\pi)^{1/2}}\Big[\mf{A}^{(1)a}_{1i}(k^{+},\v z)+\mf{A}^{(3)a}_{1i}(k^{+},\v z)\Big]\ket{0}
+\mf{N}\ket{g^{a}_{i}(k^{+},\v z)}\notag\\
&+\frac{1}{2\pi}\int dp^{+}\int d^{2}x\,\mf{B}^{ba}_{3ji}(k^{+},\v z;p^{+},-\v x)\ket{g^{b}_{j}(p^{+},\v x)}\notag\\
&+\frac{1}{(2\pi)^{1/2}}\int dp^{+}\int d^{2}x\,\Big[\mf{A}^{(3)b}_{2j}(p^{+},-\v x)-\mf{A}^{(1)\dagger b}_{1j}(p^{+},\v x)\Big]
\ket{g^{a}_{i}(k^{+},\v z)\,g^{b}_{j}(p^{+},\v x)}\notag\\
&+\frac{1}{(2\pi)^{1/2}}\int dp^{+}\int d^{2}x\,d^{2}y\,\mf{C}^{bca}_{1jki}(k^{+},\v z-\v y;p^{+},\v y-\v x)
\ket{g^{b}_{j}(p^{+},\v x)\,g^{c}_{k}(k^{+}-p^{+},\v y)}+\dots
\end{align}
In terms of the longitudinal momentum fraction $\xi=p^{+}/k^{+}$,
\begin{align}\label{Omegaongposzeta}
\Omega\ket{g^{a}_{i}(k^{+},\v z)} =\ &\frac{1}{(2\pi)^{1/2}}\Big[\mf{A}^{(1)a}_{1i}(k^{+},\v z)+\mf{A}^{(3)a}_{1i}(k^{+},\v z)\Big]\ket{0}
+\mf{N}\ket{g^{a}_{i}(k^{+},\v z)}\notag\\
&+\frac{k^{+}}{2\pi}\int d\xi\int d^{2}x\,\mf{B}^{ba}_{3ji}(k^{+},\v z;\xi k^{+},-\v x)\ket{g^{b}_{j}(\xi k^{+},\v x)}\notag\\
&+\frac{k^{+}}{(2\pi)^{1/2}}\int d\xi\int d^{2}x\,\Big[\mf{A}^{(3)b}_{2j}(\xi k^{+},-\v x)-\mf{A}^{(1)\dagger b}_{1j}(\xi k^{+},\v x)\Big]
\ket{g^{a}_{i}(k^{+},\v z)\,g^{b}_{j}(\xi k^{+},\v x)}\notag\\
&+\frac{k^{+}}{(2\pi)^{1/2}}\int d\xi\int d^{2}x\,d^{2}y\,\mf{C}^{bca}_{1jki}(k^{+},\v z-\v y;\xi k^{+},\v y-\v x)
\ket{g^{b}_{j}(\xi k^{+},\v x)\,g^{c}_{k}((1-\xi)k^{+},\v y)}+\dots
\end{align}
In the $\mf{B}_{3}$ and $\mf{A}_{2}$ terms $\Lambda/k^{+}<\xi<\vee/k^{+}$. In the $\mf{C}_{1}$ term, where both
gluons are softer than the incoming one, $\Lambda/k^{+}<\xi<1-\Lambda/k^{+}$.

\subsection{Light-cone wave functions}\label{subsec:lcwf}
In this subsection we collect the light-cone wave functions (LCWFs) that fix the coefficients of
Eq.~\eqref{eq:OmegaConstrained} to the order needed for the single inclusive gluon cross section at
next-to-leading order. For the soft vacuum as initial state they are known~\cite{Lublinsky:2016meo}. For the
one-gluon initial state we compute them in LCPT, component by component. They fix $\mf{B}_{3}$ and
$\mf{C}_{1}$, and they give $\mf{A}_{1}$ independently of the unitarity relations. In the diagrams of this
subsection the straight horizontal line represents the valence charges, curly lines are soft gluons, and
a vertical dashed line marks an intermediate state whose energy enters a denominator. A component of the
dressed one-gluon state is denoted by $\ket{\Psi^{a}_{i,\,X}(k)}$. Here $(k,a,i)$ are the momentum, colour
and polarisation of the incoming gluon, and the subscript $X$ lists the outgoing soft gluons and the
valence charge densities, as for the vacuum components below. The momenta of the outgoing and intermediate
gluons are integrated over, so each component depends only on the quantum numbers of the incoming gluon.
In all wave functions the state stands directly after the integration measure and is followed by its
amplitude, and we use $\int_{\v x,\v z}\equiv\int d^{2}x\,d^{2}z$, etc.

\subsubsection{The vacuum incoming state}\label{sec:vacin}
The dressing of the soft vacuum, $\ket{\Psi}=\Omega\ket{0}$, was computed to order $g^{3}$
in~\cite{Lublinsky:2016meo}, including quarks. We do not repeat the derivation and quote only the results
we need, for pure Yang--Mills theory ($N_{f}=0$). The dressed vacuum is
\begin{align}\label{NLOWF}
\ket{\Psi}=\mf{N}\ket{0}+\ket{\Psi^{\rm LO}_{g\rho}}+\ket{\Psi_{gg\rho}}+\ket{\Psi_{gg\rho\rho}}
+\ket{\Psi_{g\rho}}+\ket{\Psi_{g\rho\rho}}+\ket{\Psi_{g\rho\rho\rho}}+\dots,
\end{align}
where the subscripts list the soft gluons and the number of valence charge densities. The component
$\ket{\Psi^{\rm LO}_{g\rho}}$ is of order $g$, the two-gluon components are of order $g^{2}$ and the
remaining one-gluon components are of order $g^{3}$.

\paragraph{One gluon at leading order.}
\begin{align}\label{lowav1mom}
\ket{\Psi^{\rm LO}_{g\rho}} =-\int_{\Lambda}^{\vee}dk^{+}\int d^{2}\v k\,\ket{g^{a}_{i}(k^{+},\v k)}\,
\frac{g\,\v k^{i}}{2\pi^{3/2}\sqrt{k^{+}}\,\v k^{2}}\,\rho^{a}(-\v k)\,.
\end{align}
In transverse coordinates,
\begin{align}\label{lowav1}
\ket{\Psi^{\rm LO}_{g\rho}} =\frac{ig}{2\pi^{3/2}}\int_{\Lambda}^{\vee}\frac{dk^{+}}{\sqrt{k^{+}}}\int_{\v x,\v z}
\ket{g^{a}_{i}(k^{+},\v z)}\,\frac{(x-z)^{i}}{(x-z)^{2}}\,\rho^{a}(\v x)\,.
\end{align}

\paragraph{Two gluons.}
\begin{align}\label{psiggrhomomre}
\ket{\Psi_{gg\rho}} =\ &-\int_{\Lambda}^{\vee}dk^{+}\int_{\Lambda/k^{+}}^{1-\Lambda/k^{+}}d\xi\int d^{2}\v k\,d^{2}\v p\,
\ket{g^{b}_{j}\big((1-\xi)k^{+},(1-\xi)\v k-\v p\big)\,g^{a}_{i}\big(\xi k^{+},\xi\v k+\v p\big)}\notag\\
&\times\frac{ig^{2}f^{abc}\,\rho^{c}(-\v k)\,\sqrt{\xi(1-\xi)}}{8\pi^{3}\big(\xi(1-\xi)\v k^{2}+\v p^{2}\big)\v k^{2}}
\bigg[\frac{\delta^{ij}}{2}\big((1-2\xi)\v k^{2}-2\,\v k\cdot\v p\big)+\frac{\v k^{j}\v p^{i}}{\xi}+\frac{\v k^{i}\v p^{j}}{1-\xi}\notag\\
&\hspace{1cm}+\frac{\v k^{2}\,(\xi\v k^{i}+\v p^{i})\big((1-\xi)\v k^{j}-\v p^{j}\big)
\big[\xi\big((1-\xi)\v k-\v p\big)^{2}-(1-\xi)(\xi\v k+\v p)^{2}\big]}
{2\xi(1-\xi)\,(\xi\v k+\v p)^{2}\,\big((1-\xi)\v k-\v p\big)^{2}}\bigg]\,.
\end{align}
In transverse coordinates, with the integration variable relabelled, $\xi\to1-\xi$, so that the gluon at
$\v z$ carries the fraction $\xi$,
\begin{align}\label{ggrhopos}
\ket{\Psi_{gg\rho}} =\ &\int_{\Lambda}^{\vee}dk^{+}\int_{\Lambda/k^{+}}^{1-\Lambda/k^{+}}d\xi\int_{\v x,\v z,\v w}
\ket{g^{b}_{j}(\xi k^{+},\v z)\,g^{a}_{i}((1-\xi)k^{+},\v w)}\,
\frac{ig^{2}f^{abc}\,\rho^{c}(\v x)\,\sqrt{\xi(1-\xi)}}{8\pi^{3}\big[\xi(x-z)^{2}+(1-\xi)(x-w)^{2}\big]}\notag\\
&\times\bigg[\frac{\delta^{ij}\big((x-z)^{2}-(x-w)^{2}\big)}{2(z-w)^{2}}+\frac{(x-w)^{i}(z-w)^{j}}{\xi\,(z-w)^{2}}
+\frac{(x-z)^{j}(z-w)^{i}}{(1-\xi)(z-w)^{2}}\notag\\
&\hspace{1cm}+\frac{(x-z)^{j}(x-w)^{i}}{2\xi\,(x-z)^{2}}-\frac{(x-z)^{j}(x-w)^{i}}{2(1-\xi)(x-w)^{2}}\bigg]\,.
\end{align}
\begin{align}\label{ggrhorhomom}
\ket{\Psi_{gg\rho\rho}} =\ &\int_{\Lambda}^{\vee}dk^{+}\int_{\Lambda/k^{+}}^{1-\Lambda/k^{+}}d\xi\int d^{2}\v k\,d^{2}\v p\,
\ket{g^{b}_{j}\big((1-\xi)k^{+},\v k-\v p\big)\,g^{a}_{i}(\xi k^{+},\v p)}\notag\\
&\times\frac{g^{2}\,\v p^{i}\,(\v k^{j}-\v p^{j})}{16\pi^{3}\sqrt{\xi(1-\xi)}\;\v p^{2}\,(\v k-\v p)^{2}}\,
\big\{\rho^{a}(-\v p),\rho^{b}(-\v k+\v p)\big\}\,.
\end{align}
In transverse coordinates,
\begin{align}\label{ggrhorhopos}
\ket{\Psi_{gg\rho\rho}} =\ &-\int_{\Lambda}^{\vee}dk^{+}\int_{\Lambda/k^{+}}^{1-\Lambda/k^{+}}d\xi\int_{\v x,\v y,\v z,\v w}
\ket{g^{b}_{j}((1-\xi)k^{+},\v w)\,g^{a}_{i}(\xi k^{+},\v z)}\notag\\
&\times\frac{g^{2}\,(y-z)^{i}\,(x-w)^{j}}{16\pi^{3}\sqrt{\xi(1-\xi)}\;(y-z)^{2}\,(x-w)^{2}}\,
\big\{\rho^{a}(\v y),\rho^{b}(\v x)\big\}\,.
\end{align}

\paragraph{One gluon at order $g^{3}$.}
The component linear in $\rho$ contains the one-loop corrections. With $\mu$ the $\overline{\rm MS}$ scale,
\begin{align}\label{combinedgrho}
\ket{\Psi_{g\rho}} =\ &\int_{\Lambda}^{\vee}\frac{dk^{+}}{\sqrt{k^{+}}}\int d^{2}\v k\,\ket{g^{a}_{i}(k^{+},\v k)}\,
\frac{g^{3}N_{c}\,\v k^{i}}{32\pi^{7/2}\,\v k^{2}}\,\rho^{a}(-\v k)
\bigg\{\Big[-\frac{2}{\epsilon}+\ln\frac{\v k^{2}}{\mu^{2}}\Big]
\Big[\frac{11}{3}+\ln\frac{\Lambda}{k^{+}}-2\ln\frac{\vee}{\Lambda}\Big]\notag\\
&\hspace{6cm}+\frac{\pi^{2}}{3}-\frac{67}{9}+\ln^{2}\frac{\Lambda}{k^{+}}-\ln^{2}\frac{\vee}{k^{+}}\bigg\}\notag\\
&+\int_{\Lambda}^{\vee}\frac{dk^{+}}{\sqrt{k^{+}}}\int_{\Lambda/k^{+}}^{1-\Lambda/k^{+}}d\xi\int d^{2}\v k\,
\ket{g^{a}_{i}(\xi k^{+},\v k)}\,\frac{g^{3}N_{c}\,\v k^{i}}{32\pi^{7/2}\sqrt{\xi}\,(1-\xi)\,\v k^{2}}\,\rho^{a}(-\v k)\notag\\
&\hspace{1cm}\times\bigg\{(\xi-2)\ln\xi+(\xi+1)\ln(1-\xi)-(\xi-2)\Big[-\frac{2}{\epsilon}+\ln\frac{\v k^{2}}{\mu^{2}}\Big]\bigg\}\,.
\end{align}
In transverse coordinates we combine the two terms. In the second we use the longitudinal momentum of the
emitted gluon, $l^{+}=\xi k^{+}$, as the integration variable and carry out the $\xi$ integral,
neglecting terms suppressed by powers of $\Lambda/l^{+}$. We then rename $l^{+}\to k^{+}$. The Fourier
transform of $\ln(\v k^{2}/\mu^{2})$ gives $-2\gamma_{E}-\ln\big((x-z)^{2}\mu^{2}/4\big)$, and
\begin{align}\label{grhoposition}
\ket{\Psi_{g\rho}} =\ &-\frac{ig^{3}N_{c}}{32\pi^{7/2}}\int_{\Lambda}^{\vee}\frac{dk^{+}}{\sqrt{k^{+}}}\int_{\v x,\v z}
\ket{g^{a}_{i}(k^{+},\v z)}\,\frac{(x-z)^{i}}{(x-z)^{2}}\,\rho^{a}(\v x)\notag\\
&\times\bigg\{\Big[-\frac{2}{\epsilon}-2\gamma_{E}-\ln\frac{(x-z)^{2}\mu^{2}}{4}\Big]
\Big[\frac{11}{3}+2\ln\frac{\Lambda}{k^{+}}+\ln\Big(1-\frac{k^{+}}{\vee}\Big)\Big]\notag\\
&\hspace{1cm}+\frac{\pi^{2}}{3}-\frac{67}{9}+\ln^{2}\Big(1-\frac{k^{+}}{\vee}\Big)
+\ln\Big(1-\frac{k^{+}}{\vee}\Big)\ln\frac{\vee}{k^{+}}\bigg\}\,.
\end{align}
The component quadratic in $\rho$ is
\begin{align}\label{grhorhomom}
\ket{\Psi_{g\rho\rho}} =\int_{\Lambda}^{\vee}\frac{dk^{+}}{\sqrt{k^{+}}}\int d^{2}\v k\,d^{2}\v p\,\ket{g^{a}_{i}(k^{+},\v k)}\,
\frac{ig^{3}f^{abc}}{32\pi^{9/2}\,\v p^{2}\,\v k^{2}}\,\big\{\rho^{b}(-\v p),\rho^{c}(-\v k+\v p)\big\}\,
\mathcal{K}^{i}(\v k,\v p;k^{+})\,,
\end{align}
with
\begin{align}\label{grhorhokernel}
\mathcal{K}^{i}(\v k,\v p;k^{+})=\ &\ln\frac{\Lambda}{k^{+}}\,\frac{\v k^{2}\,\v p^{i}-(\v k\cdot\v p)\,\v k^{i}}{(\v k-\v p)^{2}}
-2\,\v k^{i}\ln\frac{\vee}{\Lambda}-2\,\v k^{i}\ln\frac{\Lambda\,\v k^{2}+k^{+}\v p^{2}}{\vee\,\v k^{2}+k^{+}\v p^{2}}\notag\\
&+\frac{(2\v p^{2}-2\,\v k\cdot\v p+\v k^{2})\,\v k^{i}+\v k^{2}\,\v p^{i}}{(\v k-\v p)^{2}}\,
\ln\Big[1+\Big(\frac{\vee}{k^{+}}-1\Big)\frac{\v k^{2}}{(\v k-\v p)^{2}}\Big]\notag\\
&-\frac{\v p^{2}\,\v k^{i}+\v k^{2}\,\v p^{i}}{(\v k-\v p)^{2}}\,\ln\frac{\vee-k^{+}}{\Lambda}-\v k^{i}\ln\frac{\vee}{k^{+}}\,.
\end{align}
The term $-2\v k^{i}\ln(\vee/\Lambda)$ has the same colour and momentum structure as the
commutator~\eqref{eq:NAcomm}. In transverse coordinates,
\begin{align}\label{grhorhopos}
\ket{\Psi_{g\rho\rho}} =\ &\frac{1}{2\pi}\int_{\Lambda}^{\vee}\frac{dk^{+}}{\sqrt{k^{+}}}\int_{\v x,\v y,\v z}
\ket{g^{a}_{i}(k^{+},\v z)}\,\frac{ig^{3}f^{abc}}{32\pi^{9/2}}\,\big\{\rho^{b}(\v x),\rho^{c}(\v y)\big\}\notag\\
&\times\int d^{2}\v k\,d^{2}\v p\;\frac{e^{i\v k\cdot(\v z-\v y)}\,e^{i\v p\cdot(\v y-\v x)}}{\v p^{2}\,\v k^{2}}\,
\mathcal{K}^{i}(\v k,\v p;k^{+})\,.
\end{align}
The component cubic in $\rho$ is
\begin{align}\label{3rhosmom}
\ket{\Psi_{g\rho\rho\rho}} =\frac{g^{3}}{16\pi^{9/2}}\ln\Big(\frac{\vee}{\Lambda}\Big)\int_{\Lambda}^{\vee}\frac{dk^{+}}{\sqrt{k^{+}}}
\int d^{2}\v k\,d^{2}\v p\,\ket{g^{a}_{i}(k^{+},\v k)}\,\frac{\v k^{i}}{\v k^{2}\,\v p^{2}}\,
\rho^{b}(\v p)\,\rho^{b}(-\v p)\,\rho^{a}(-\v k)\,,
\end{align}
and in transverse coordinates
\begin{align}\label{3rhospos}
\ket{\Psi_{g\rho\rho\rho}} =\ &-\frac{ig^{3}}{16\pi^{9/2}}\ln\Big(\frac{\vee}{\Lambda}\Big)\int_{\Lambda}^{\vee}\frac{dk^{+}}{\sqrt{k^{+}}}
\int_{\v w,\v x,\v y,\v u,\v z}\ket{g^{a}_{i}(k^{+},\v z)}\notag\\
&\times\frac{(x-u)\cdot(y-u)\,(w-z)^{i}}{(x-u)^{2}\,(y-u)^{2}\,(w-z)^{2}}\,\rho^{b}(\v x)\,\rho^{b}(\v y)\,\rho^{a}(\v w)\,.
\end{align}
Comparing with Eqs.~\eqref{lowav1mom} and~\eqref{eq:Nmom}, $\ket{\Psi_{g\rho\rho\rho}}=\mf{N}^{(2)}\ket{\Psi^{\rm LO}_{g\rho}}$. It
comes from the canonical normalisation of the third-order state, and $\mf{N}^{(2)}$ stands to the left of
$\rho^{a}$. This ordering matters, since $\mf{N}^{(2)}$ does not commute with $\rho^{a}$
[cf.\ Eq.~\eqref{eq:NAcomm}].

\paragraph{Normalisation.}
The condition $\braket{\Psi|\Psi}=1$ fixes $\mf{N}$:
\begin{align}\label{eq:Ncondition}
\mf{N} =\ &1-\frac{1}{2}\Big[\braket{\Psi^{\rm LO}_{g\rho}|\Psi^{\rm LO}_{g\rho}}
+\braket{\Psi^{\rm LO}_{g\rho}|\Psi_{g\rho}}+\braket{\Psi_{g\rho}|\Psi^{\rm LO}_{g\rho}}
+\braket{\Psi^{\rm LO}_{g\rho}|\Psi_{g\rho\rho}}+\braket{\Psi_{g\rho\rho}|\Psi^{\rm LO}_{g\rho}}\notag\\
&\hspace{1cm}+\braket{\Psi^{\rm LO}_{g\rho}|\Psi_{g\rho\rho\rho}}+\braket{\Psi_{g\rho\rho\rho}|\Psi^{\rm LO}_{g\rho}}
+\braket{\Psi_{gg\rho}|\Psi_{gg\rho}}+\braket{\Psi_{gg\rho}|\Psi_{gg\rho\rho}}\notag\\
&\hspace{1cm}+\braket{\Psi_{gg\rho\rho}|\Psi_{gg\rho}}+\braket{\Psi_{gg\rho\rho}|\Psi_{gg\rho\rho}}\Big]
-\frac{1}{8}\braket{\Psi^{\rm LO}_{g\rho}|\Psi^{\rm LO}_{g\rho}}\braket{\Psi^{\rm LO}_{g\rho}|\Psi^{\rm LO}_{g\rho}}\,,
\end{align}
where the overlaps are operators on the valence space. Only the first term is of order $g^{2}$. All the
others are of order $g^{4}$: $\mf{N}$ enters the cross section multiplied by terms of order $g^{2}$, so
they contribute only at order $g^{6}$. To the accuracy needed,
\begin{align}\label{eq:Nmom}
\mf{N}=1+\mf{N}^{(2)}
=1-\frac{g^{2}}{8\pi^{3}}\int_{\Lambda}^{\vee}\frac{dp^{+}}{p^{+}}\int d^{2}\v p\,\frac{\rho^{a}(-\v p)\,\rho^{a}(\v p)}{\v p^{2}}
=1-\frac{g^{2}}{8\pi^{3}}\ln\Big(\frac{\vee}{\Lambda}\Big)\int d^{2}\v p\,\frac{\rho^{a}(-\v p)\,\rho^{a}(\v p)}{\v p^{2}}\,,
\end{align}
in agreement with Eq.~\eqref{eq:Nrel} and with $\mf{N}^{(2)}$ of Eq.~\eqref{eq:A1N2explicit}. In transverse
coordinates,
\begin{align}\label{eq:Npos}
\mf{N}=1-\frac{g^{2}}{8\pi^{3}}\ln\Big(\frac{\vee}{\Lambda}\Big)\int_{\v x,\v y,\v z}
\frac{X\cdot Y}{X^{2}\,Y^{2}}\,\rho^{a}(\v x)\,\rho^{a}(\v y)\,,\qquad X\equiv\v x-\v z\,,\quad Y\equiv\v y-\v z\,.
\end{align}
$\mf{N}^{(2)}$ is a non-local operator on the valence space: it correlates charges at different transverse
points through the kernel $X\cdot Y/(X^{2}Y^{2})$. This is why it does not commute with $\rho^{a}(\v x)$,
Eqs.~\eqref{eq:NAcomm} and~\eqref{eq:NAcommx}.

Comparing Eq.~\eqref{NLOWF} with Eq.~\eqref{eq:Omegavac}, where
$(2\pi)^{3/2}\int d^{3}p/(2\pi)^{3}=(2\pi)^{-3/2}\int d^{3}p$, identifies the coefficients.
$(2\pi)^{-3/2}\mf{A}^{(1)}_{2}$ is the amplitude of $\ket{\Psi^{\rm LO}_{g\rho}}$, which gives
Eq.~\eqref{eq:A1N2explicit}. $(2\pi)^{-3/2}\mf{A}^{(3)}_{2}$ is the sum of the amplitudes of
$\ket{\Psi_{g\rho}}$, $\ket{\Psi_{g\rho\rho}}$ and $\ket{\Psi_{g\rho\rho\rho}}$. $(2\pi)^{-3}\mf{B}_{2}$ is the sum of
the amplitudes of $\ket{\Psi_{gg\rho}}$ and $\ket{\Psi_{gg\rho\rho}}$, up to the Jacobian
$d^{3}p\,d^{3}q=k^{+}dk^{+}d\xi\,d^{2}\v k\,d^{2}\v p$ of the change of variables.
